\documentclass[11pt]{article}
\usepackage[margin=.9in]{geometry}
\usepackage[T1]{fontenc}
\usepackage[colorlinks=true,urlcolor=blue]{hyperref}
\newcommand{\entry}[3]{\par\smallskip\noindent\textbf{#1. #2.} #3\par}
\title{Swarm Coordination Fault Inventory\\Cleanroom Synthesis}
\author{Working draft}
\date{24 November 2026}
\begin{document}
\maketitle

\section{What the inventory classifies}
The unit is a service obligation: who owes what to whom, under which conditions, by when, and what counts as evidence of fulfillment. An obligation can be local, such as delivery of a fresh estimate, or joint, such as keeping a shared load supported. The team includes any infrastructure, human supervisor or external supplier on which the work depends. Authority can be centralized, distributed or mixed.

A \emph{fault mechanism} is a cause or condition that can break a required service. A \emph{breach} is a departure from the applicable obligation. Record them separately: one fault can break several services, a latent fault need not have broken any, and correct components can collectively breach a joint obligation. A closure or lawful withdrawal is not by itself a fault. The obligation it triggers might be transfer, rerouting, degraded service, or a protected stop. Separately record deficiencies in the operating rules or in the research model. If the written rules permit harm they were supposed to prevent, name the stakeholder and the protection missing from the rules; do not report an invented breach of the written contract.

Two workers independently produced the first inventories using only their own briefs and model knowledge, without repository, incident or prior conversation material. This synthesis reorganizes their returns. The comparison in the last section is retrospective and is not part of the cleanroom. The two runs are independent in supplied context, not independent sources of empirical evidence. Illustrations are hypothetical; the inventory does not establish frequencies, exhaustive coverage or the feasibility of any response.

\section{Obligations and distinguishing observations}

\subsection{Observing and sharing the situation}
\entry{S01}{Coverage}{An event, entity or region cannot be observed within the promised interval or envelope. Does the consumer know the difference between unseen and clear?}
\entry{S02}{Validity and attribution}{A measurement, entity association or confidence claim is wrong at its stated time and place. Compare with an independent reference or two histories that disambiguate the source.}
\entry{S03}{Comparable state}{Locally correct observations become inconsistent when combined across clocks, coordinate frames, units, versions, feature identifiers or uncertainty conventions. Trace the conversions used by each consumer.}
\entry{S04}{Transfer and freshness}{A required update is omitted, late, duplicated, corrupted, reordered, delivered asymmetrically or suppressed behind a queue. Compare what each recipient actually had by its decision deadline, not the sender's submission receipt.}
\entry{S05}{Consumption and interpretation}{Correctly received information is applied to the wrong entity, epoch or decision, misread, or used after expiry. Determine whether the receiver used the intended meaning and distinguished unknown from clear.}
\entry{S06}{Reconciliation of shared state}{Conflicting valid-looking replicas, reservations or logs remain after a merge. Check whether a concurrent update, old release or replay applies to the actual object and generation.}

\subsection{Defining and assigning the work}
\entry{S07}{Operating intent and obligations to outsiders}{A known constraint is omitted, precedences diverge, or a proxy replaces what operators, third parties or teammates were owed. A snow-clearing team can meet a road-clearance target while blocking required pedestrian access. Name the actual obligation and whether the published rule encoded it.}
\entry{S08}{Membership, role and effective capability}{The wrong body is counted or assigned a role, a returning member is treated as continuously present, or readiness is asserted on capability no longer available. Trace participant, role and effective capability at the decision time.}
\entry{S09}{Authority and arbitration}{Conflicting commands are not resolved under the governing precedence, authorization is absent or obsolete at effect time, or a competent authority cannot intervene when required. Presence, technical control and lawful authority are different facts.}
\entry{S10}{Compatibility of shared commitments}{Each participant holds a locally meaningful plan or reservation, but the commitments cannot hold together over a named interval. Different views alone do not establish incompatibility.}
\entry{S11}{Allocation, eligibility and custody}{Required work is unowned, assigned more than once where exclusive, or handed off before the successor accepts and can act. Trace the object and its actual support or custody as well as the assignment record.}
\entry{S12}{External change enters the decision}{A relevant external condition is correctly available and understood but not used by the decision policy. A correctly represented closure ignored by planning differs from an unseen closure or a restriction that did not apply.}

\subsection{Coordinating physical action}
\entry{S13}{Prerequisites and phase}{A dependent action begins before its prerequisite, after its window or without required overlap. Trace actual effects and evidence times separately. In a handoff of a shared load, receipt of acceptance need not establish that the other carrier is supporting it.}
\entry{S14}{Joint physical and resource compatibility}{Locally feasible actions create incompatible trajectories, loads, force profiles, timing or concurrent resource demand. Two carriers can keep within their individual motion limits while their combined motion damages a flexible load. Superpose actions at the shared interface.}
\entry{S15}{Supply, margins and endurance}{Available power, time, compute, braking, payload or support capacity drops below what is needed for the full obligation. Check margins and replenishment across the promised horizon, not merely readiness for the next action.}
\entry{S16}{Correlated dependencies}{Apparently independent vehicles or evidence sources share a physical resource, map, correction feed, power or judgment model. Simultaneous failure does not establish a common cause; trace the shared dependency.}
\entry{S17}{Physical realization}{A command is omitted, stuck, delayed, over-applied, or produces the wrong effect under changed dynamics. Distinguish bad commands faithfully followed, motion achieved without task effect, and task effect achieved without confirmation.}
\entry{S18}{Collective progress}{The team stalls, starves a member, cycles assignments or oscillates despite local activity. Identify the missed progress obligation; waiting can be a compliant protected response or a service breach depending on what remained owed.}
\entry{S19}{Transition and mixed modes}{Before and after modes are separately acceptable but their overlap or gap is not. Check intermediate states, including one member entering return mode while another still uses the nominal route.}

\subsection{Supervision, recovery and accounting}
\entry{S20}{Human and organizational supervision}{A dashboard conflates accepted stop with standstill, an alarm has no responsible recipient, or a human and automation both act or both wait through a takeover. Ask separately whether the supervisor received enough evidence in time and had effective authority to act.}
\entry{S21}{Detection, escalation and containment}{A breach is missed, attributed to the wrong cause or escalated after viable options expire; locally defensible isolation or fallback creates a further joint loss. State which actors know which facts at effect time.}
\entry{S22}{Restoration and re-entry}{Restored connectivity or a repaired component does not validate commitments, prior effects or membership. Verify the preconditions for re-entry; residual uncertainty can rule out repeating an irreversible action even when an inspection remains admissible.}
\entry{S23}{Effects and completion accounting}{A claim of completion, discharge or resource consumption does not match work actually done. Verify partial, duplicate, omitted or reversed effects and effects for which no confirmation survived.}
\entry{S24}{Protection after apparent closure}{Cancellation, completion or local stopping leaves a bystander hazard, blocked access, outstanding duty or released object. Check the actual terminal state and obligations that persist.}
\entry{S25}{Evidence for later review}{Records are insufficient to reconstruct the relevant actors, authority, information and effects. A plausible story is not a finding when competing accounts fit what was recorded.}

\section{Mechanism tags and interaction tests}
Record origin, scope, actual dependency, temporal course and operating phase separately from the obligation ID. The same defect can be intermittent, shared and visible only during transition without being three independent faults. A practical record identifies the change and comparison case; onset, when each actor can detect it, and when it becomes actionable; effects and duties that survive; and the evidence or missing observation for each claimed breach. Consequence is useful for triage; source and dependency are useful for assigning repair ownership. Neither can replace identification of what failed.

\entry{I01}{Threshold and accumulation}{Several small losses jointly exhaust a required reserve or close a timing window. Compare what the group can still do, not merely what modes it reports.}
\entry{I02}{Complementary loss of evidence}{A coverage failure removes one source and stale information makes the backup unusable. Each change alone leaves a means to establish the prerequisite.}
\entry{I03}{Common cause versus coincidence}{Compare a shared dependency failure to separate losses with the same component-level observations. The distinction changes what counts as replacement and what justifies confidence.}
\entry{I04}{Mediation through obligations}{A missed update yields incompatible reservations, which in turn produce capacity conflict. Check the intermediate change and its consumer rather than inferring the chain from start and end events alone.}
\entry{I05}{Masking and genuine compensation}{A stoppage prevents a space conflict while stale authority remains live. Compare an isolated symptom suppression with a recovery that actually discharges the obligation; the two may diverge on resumption.}
\entry{I06}{Order, overlap and persistence}{A role change immediately before recovery can differ from one after recovery. Vary event order and what persists after initiators clear.}
\entry{I07}{An adaptive response sustains the fault}{Repeated reassignment consumes the shared capacity needed to reconcile assignments. A periodic symptom alone does not establish this loop; trace the response's effect on the next round.}

Where practical, compare baseline, each individual change and the combination on the same relevant exogenous choices. Cross event order as well as simultaneous onset. Separate observed co-occurrence from supported interaction; replay claims inherit the bounds of the replay model.

\section{How to look for omitted classes}
Pick an operation distant from the system used to write the list. Follow one coupled obligation from commissioning to closure and revisit it at admission, reassignment, withdrawal, mixed-mode transition, pause, restart and exhaustion. Use at least two substantially different services: for example, delivery through a warehouse's shared passage and continuity of support for a marine survey. Check loads and custody in a cooperative transport example; check obligations to bystanders or third parties in a public-space example. These are construction prompts, not selected simulations.

At every boundary test omission, excess or duplication, wrong referent, mistiming, divergence between consumers, concealed loss and incomplete restoration. Include external suppliers and people even when they issue no motion commands. Seek both counterexamples: every local check passes but a joint obligation fails, and the visible joint output looks correct but authority, protection or the evidence needed to verify it fails. A new class needs a named duty, a comparison and a witness for an observably or operationally different failure. When only the platform changes, add an illustration. When no available observation can distinguish the alternatives, preserve an unresolved finding with the evidence that would decide it. A fault-free but unnecessarily costly permitted response belongs in an efficiency review, not a safety tally.

\section{Comparison with the earlier draft}
The \href{https://chatgpt.com/space/page_67998e48ab148191b64fd6759e583bff}{earlier inventory} already treats mixed revisions, correlated observations, receipt versus execution, residual occupancy and live obligations after recovery. Their reappearance without that draft is a useful check on the classification. It does not validate them empirically. Four differences are worth carrying forward:
\begin{description}
\item[Supervision and takeover (S20).] The earlier F26/F45/F53 rows cover authority, alarms and support separately. A dedicated supervision boundary makes the joint obligation visible: who can intervene, who is expected to, and what view supports the decision. An operator with lawful authority but no reliable indication of standstill is not the same failure as an operator with truthful status but no effective control.
\item[Rules can be wrong (S07).] Earlier F32 treats inconsistent application of a legitimate rule change and A02 treats checking a proxy. The blind returns explicitly add a different obligation: the operative rule can be faithfully applied and itself allow a result that harms somebody it was meant to protect.
\item[A fact can be known and still omitted (S12).] Earlier F06/F14/F25 provide related entries. The explicit split yields a better audit: compare late receipt of a closure, incorrect interpretation of a timely receipt, and a correctly understood closure omitted by planning. Observing the same crossing does not identify which repair is required.
\item[An apparent repair can hide a live fault (I05).] Earlier F48 records residual effects. The blind work adds a testable interaction: the second event can suppress evidence of the first without repairing it. Compare the same stopped state with and without stale authority at re-entry. No extra headline class is needed if this distinction is preserved in the test register.
\end{description}
Conversely, the earlier draft is sharper about losing a communication modality's semantics (F12), an observation action changing geometry and available service (F13), and a policy whose premises expire before installation (F31). The blind classes can accommodate these, but did not specify them at comparable resolution. Use the cleanroom obligation structure as the core, retain these mechanism-level refinements, and use incidents later to test coverage rather than define the classes. Neither contract v0.5 nor the earlier draft was changed.
\end{document}
